% SPDX-FileCopyrightText: 2026 Hasan Melih Akbulut
% SPDX-License-Identifier: CC-BY-4.0
\documentclass[11pt,a4paper]{article}

\usepackage[T1]{fontenc}
\usepackage[utf8]{inputenc}
\usepackage{lmodern}
\usepackage[margin=1in]{geometry}
\usepackage{booktabs}
\usepackage{array}
\usepackage{amsmath}
\usepackage{siunitx}
\usepackage{microtype}
\usepackage{tabularx}
\usepackage[hidelinks]{hyperref}

\sisetup{group-separator={,}, group-minimum-digits=4,
         group-digits=integer, text-series-to-math=true}

\newcolumntype{L}[1]{>{\raggedright\arraybackslash}p{#1}}
\newcommand{\file}[1]{\texttt{\small #1}}

\title{A Fault-Tolerant Neuromorphic System-on-Chip\\
in an Open \SI{130}{\nano\meter} PDK:\\
\large Architecture, Protection Cost and Physical Implementation}

\author{Hasan Melih Akbulut\thanks{Independent researcher.
Correspondence: \texttt{57303760+Melihakbulut221@users.noreply.github.com}.
Sources, harnesses and the full design record are open; see
\S\ref{sec:artefacts}.}}

\date{14 September 2026}

\begin{document}
\maketitle

\begin{abstract}
We describe a system-on-chip that places a RISC-V management core and an
event-driven spiking-neural-network node on one die in the IHP SG13G2
open \SI{130}{\nano\meter} process, hardened through an entirely open
flow, and we report what its fault-tolerance costs where the cost is
measurable.

The design is \num{168737} placed instances around eight vendor SRAM
macros on a \SI{4.79}{\milli\meter\squared} die. Protection is layered
rather than uniform: the register file carries a \num{39}/\num{32}
SECDED code with a background scrub, the system memory carries four
byte-wide \num{16}/\num{8} codewords per \num{64}-bit row so that a byte
store needs no read-modify-write, the boot ROM carries its own
\num{39}/\num{32} word code, the watchdog's state is triplicated, and the
accelerator's \num{55} configuration bits are triplicated in the netlist.
Every correction and every uncorrectable syndrome is counted in a
telemetry block that lives in the power-on domain and survives the reset
the watchdog causes.

\textbf{The costs are stated as measured pairs, not as estimates.}
Protecting the system memory costs \num{2068} cells, \num{234}
flip-flops and \SI{31391}{\micro\meter\squared}---\SI{4.78}{\percent} of
the standard-cell area---and halves the usable RAM from \num{64} to
\num{32}\,KiB, because the protection buys its check bits out of the
word width rather than out of a seventh macro. Protecting the register
file costs \num{222} flip-flops.

The layout routes with zero detailed-routing violations, zero
disconnected pins and zero antenna violations, matches its netlist
uniquely under Netgen, and returns zero XOR differences. It does
\emph{not} meet its timing constraint: \SI{-7.35}{\nano\second} of setup
slack at the slow corner against a \SI{20}{\nano\second} period, on
\num{3550} violating endpoints. We report that with the same weight as
the results that pass, and \S\ref{sec:limits} states the four remedies
that were measured and what each returned.

\textbf{No silicon exists.} Nothing here has been fabricated or exposed
to a beam, and neither open PDK involved publishes total-ionising-dose or
single-event data. Every reliability figure is a seeded fault-injection
measurement on register-transfer or gate-level models.
\end{abstract}

\section{Introduction}
\label{sec:intro}

Spacecraft that want to run inference on sparse, event-driven sensor data
have recently acquired a commercial option: Frontgrade Gaisler's
GR801~\cite{gr801}
pairs a fault-tolerant NOEL-V RV64GC core with a licensed BrainChip
Akida accelerator on \SI{28}{\nano\meter} FDSOI, under a national space
agency contract. It is the reference this work measures itself against
and it is ahead on every axis a mission cares about: process node,
qualification status, flight heritage and the fact that it will exist as
silicon.

This design is not a competitor to that. It is an open one. Every
source, every harness, every measurement record and every negative
result is published, on a process design kit and a tool flow that anyone
can install. What that buys is not performance. It is that each claim
below can be re-derived by a reader rather than believed.

\paragraph{What is on the die.} A \num{32}-bit RISC-V management core
with physical memory protection; \num{72}\,KiB of on-chip memory in six
vendor macros---\num{64}\,KiB of SRAM in four and an \num{8}\,KiB boot
ROM in two---with two further macros carrying the ROM's check bits, for
eight in the layout; a system fabric with an error slave and
plug-and-play device tables; a core-local interruptor, a general-purpose
timer and a three-stage escalating watchdog; a console UART, a
\num{16}-pin GPIO port and a QSPI flash controller; and one node of an
event-driven spiking accelerator, instantiated from a frozen shuttle
submission rather than copied from it.

\paragraph{What this paper reports.} The architecture
(\S\ref{sec:arch}); the protection scheme and its measured cost
(\S\ref{sec:protect}); the verification, which is three independent
kinds and is counted (\S\ref{sec:verify}); the physical implementation
and its sign-off (\S\ref{sec:phys}); and what the design does not do,
which is a long list and is stated rather than omitted
(\S\ref{sec:limits}).

\section{Architecture}
\label{sec:arch}

\subsection{Management core}

The core is lowRISC's Ibex in its \texttt{small-pmp} configuration:
RV32IMC with a two-stage pipeline, a flip-flop register file, no
instruction cache, no branch-target ALU and no writeback stage, with
four physical-memory-protection regions at granularity~\num{0}.

It is vendored \emph{pristine}. The upstream commit is pinned, the
SystemVerilog is converted to Verilog-2005 by sv2v, and \textbf{no patch
is applied to Ibex itself}. Everything this project adds to the core is
outside it or substituted at a module boundary, which is what makes
``the core is unmodified'' a checkable statement rather than an
intention.

\texttt{SecureIbex} is \num{0}: no lockstep, no shadow register file, no
dummy instructions, no memory-interface ECC. That is a recorded decision
rather than an omission---the lockstep configuration roughly doubles the
core, and this design spends that area on memory protection instead,
where \S\ref{sec:protect} shows the upset population actually is.

\subsection{Fabric and memory map}

The system fabric speaks Ibex's own request/grant/response protocol
rather than a standard bus, with two masters (the instruction and data
ports), six slave ports, round-robin arbitration and two outstanding
requests. No-starvation is proved rather than argued.

The peripheral bus is real AMBA~3 APB---a \num{1}\,MiB window of sixteen
\num{4}\,KiB slots---and the bridge between the two is proved by
k-induction.

Two properties of the map are worth stating because they are the kind
that are usually absent. \textbf{Every unmapped address raises a bus
error} that the core takes as an access fault: there is no silent read
of zero anywhere in the map. And two plug-and-play device tables in
GRLIB's record layout let software enumerate what is present instead of
being told.

\subsection{Neuromorphic node}

The accelerator is one node of a sixteen-node address space, of which
one is instantiated and the other fifteen decode and fault. The node is
the frozen Tiny Tapeout TTIHP26b pilot, \textbf{instantiated out of its
own directory rather than copied}, with a test that fails if that
directory has been modified---so the SoC cannot quietly drift away from
the part that was submitted.

The node carries \num{8} neurons, \num{8} axons and \num{64} four-bit
signed synapses in an event-driven leaky-integrate-and-fire
organisation, with a \num{72}/\num{64} Hsiao SECDED code on the weight
word, \num{55} configuration bits triplicated in the netlist, and four
fault-report pins.

\textbf{It is driven exactly as external silicon would be.} A hardware
serial master issues one \num{40}-bit mode-0 frame per register
transaction, costing \num{176} clock cycles per register access. That is
deliberately expensive: the point is that the interface the SoC uses is
the interface a separate die would present, so the integration is not
made easy by being made unrealistic.

Events reach the node through two AER queues carrying the die's own
pointer voting, entry parity and drop counting, reported through a
\num{14}-bit interrupt cause register.

\section{Fault tolerance, and what it costs}
\label{sec:protect}

\subsection{The scheme is layered, not uniform}

Different structures get different mechanisms, chosen by what the
structure is and what an upset in it does:

\begin{center}
\begin{tabular}{L{3.4cm}L{4.3cm}L{5.2cm}}
\toprule
structure & mechanism & why this one \\
\midrule
register file & \num{39}/\num{32} SECDED, background scrub over
  \texttt{x1}--\texttt{x31} & dense, hot, and an upset propagates
  immediately into computation \\
system SRAM & four \num{16}/\num{8} byte codewords per \num{64}-bit row,
  plus a scrubber & a byte store must not become a read-modify-write \\
boot ROM & \num{39}/\num{32} word code, two check macros & its contents
  decide whether the part boots at all \\
watchdog state & triplication with a voter & small, and its failure mode
  is an unwanted reset \\
accelerator config & \num{55} bits triplicated in the netlist & written
  once, read for the lifetime of a configuration \\
\bottomrule
\end{tabular}
\end{center}

\subsection{The memory protection, and the trade it makes}

The system memory's scheme is the one with a real design consequence and
it is worth stating plainly. A \num{72}/\num{64} code would need eight
check bits that a \num{64}-bit macro row does not have, so the check bits
would need a seventh macro---costing roughly a fifth of the macro area
and putting a read-modify-write on every \texttt{sb} and \texttt{sh}.

Instead the usable word narrows: \textbf{one \num{32}-bit word per
\num{64}-bit row, as four \num{16}/\num{8} byte codewords}, each the
shortening of the one SECDED code this design already owns. A byte is
written through the macro's own bit mask with no read-modify-write, no
extra cycle and nothing for the fabric to re-prove. \textbf{The price is
that the usable RAM is \num{32}\,KiB, out of the \num{64}\,KiB the four
macros physically hold.}

That is a real cost and it is the right one for this part, whose
bring-up program's static data is \num{192} bytes---but it is a trade
that a design with a larger working set could not make, and it is stated
here as a trade rather than as a feature.

\subsection{Measured cost}

\begin{center}
\begin{tabular}{lrrr}
\toprule
& cells & flip-flops & area \\
\midrule
register-file SECDED and scrub & --- & \num{222} & --- \\
system-memory protection & \num{2068} & \num{234} &
  \SI{31391}{\micro\meter\squared} (\SI{4.78}{\percent}) \\
\bottomrule
\end{tabular}
\end{center}

Two things about that table. It is measured at the whole-SoC boundary
rather than at the block, because a block-level figure for a protection
mechanism is a number that does not survive synthesis. And the
whole-SoC cycle invariant is \emph{unmoved} at \num{217682} with the RAM
halved---the protection buys its correctness out of area and word width,
not out of time.

\subsection{Telemetry that survives the reset}

Every correction, every corrected read cycle, every uncorrectable
syndrome and every watchdog voter mismatch is counted, in four
saturating counters with four sticky status bits and a maskable
interrupt.

The design decision worth naming is where that record lives: \textbf{in
the power-on domain}, so it survives the system reset the watchdog
causes. A counter that is cleared by the event it is counting is not
telemetry, and this one is not.

\section{Verification}
\label{sec:verify}

Three independent kinds, all counted and all runnable from a checkout.

\paragraph{Formal.} \num{64} SymbiYosys tasks across \num{15} property
sets covering the fabric, the APB bridge, the memory codec, the boot
block, the CLINT, the QSPI controller, the GPIO port, the scrubber, the
watchdog and its triplicated state, and the clock-gate wake path. The
fabric's no-starvation property and the APB bridge's conformance are
both proved rather than tested.

\paragraph{Simulation.} \num{464} cocotb tests across the SoC blocks and
the pilot, and \num{520} Python tests that check the generated register
maps, the memory map, the golden model, the documentation's own counts
and the invariants this project keeps about its flow configurations.

\paragraph{Fault injection.} Seven register-transfer campaigns covering
the management core, the NPU connection, the event engine, the watchdog
and the node itself, each with per-structure outcome distributions and
stated Wilson bounds---plus a gate-level campaign that replays one of
them on the sign-off netlist, one injection for one, to see whether the
register-transfer result survives the transformation to gates.

\paragraph{And a negative result we keep.} One of those campaigns
produced \emph{bit-identical} output on a defective design and its
repaired version, which is not a pass: it is a campaign that was not
looking at the structure it was supposed to check. That finding, and
what it says about fault-injection campaigns as instruments, is the
subject of a companion paper~\cite{companion_tmr} and is named here
because a verification section that lists only its successes is a
brochure.

\section{Physical implementation}
\label{sec:phys}

\subsection{The flow}

Entirely open: Yosys for synthesis, OpenROAD through LibreLane for place
and route, Magic and KLayout for design-rule checking, Netgen for
layout-versus-schematic, OpenSTA for timing. The process kit is IHP's
open SG13G2 at a pinned commit. No step of the flow is a black box to a
reader, which is the property that makes the rest of this section
checkable.

\subsection{The layout}

\begin{center}
\begin{tabular}{lr}
\toprule
placed instances & \num{168737} \\
\quad of which standard cells & \num{64409} \\
\quad of which vendor SRAM macros & \num{8} \\
die area & \SI{4.79}{\milli\meter\squared} \\
macro area & \SI{2.34}{\milli\meter\squared} (\SI{49}{\percent} of the die) \\
routed wirelength & \SI{5.75}{\meter} \\
total power, typical corner & \SI{56.1}{\milli\watt} \\
\bottomrule
\end{tabular}
\end{center}

Synthesis of the whole SoC as one design gives \num{27694} cells,
\num{3077} flip-flops and \SI{392932.89}{\micro\meter\squared}
(\num{54.141}\,kGE), which is \SI{1.11}{\percent} \emph{smaller} than the
sum of the nine blocks synthesised separately---the whole is not the sum
of its parts, and the difference is worth knowing before anyone budgets
from block figures.

\subsection{Sign-off}

\begin{center}
\begin{tabular}{lL{8.4cm}}
\toprule
detailed-route DRC & \num{0} \\
disconnected pins & \num{0} \\
antenna violations & \num{0} nets, \num{0} pins \\
Netgen LVS & circuits match uniquely, all seven difference counters
  \num{0}, with the vendor macros abstracted to their pin lists \\
KLayout XOR & \num{0} differences \\
KLayout DRC & \num{11048} markers, \textbf{all of them inside the vendor
  SRAM macros}; \num{0} name a cell this design drew \\
Magic DRC & \num{182} boxes, \num{20} inside a macro footprint, and every
  one of the other \num{162} within \SI{0.5}{\micro\meter} of a macro
  edge \\
\bottomrule
\end{tabular}
\end{center}

The last two rows need their qualifier, and it is the subject of a
second companion paper~\cite{companion_decks}: the vendor SRAM macro does not pass this PDK's own
geometric decks, which is a property of the kit rather than of this
design, and the honest form of a deck result here is therefore two
numbers---inside the macro and outside it---rather than one.

\section{What this design does not do}
\label{sec:limits}

\paragraph{It does not meet its timing constraint.}
\SI{-7.35}{\nano\second} of setup slack at the slow corner against a
\SI{20}{\nano\second} period, on \num{3550} violating endpoints. Four
remedies were measured rather than assumed: a second floorplan, a
synthesis restructuring pass, a register on the memory read return, and
an explicit capacitance constraint. \textbf{Two of the four made it
worse}, and the diagnosis that had ranked the capacitance constraint
first modelled it at \SI{+4.3}{} to \SI{+5.5}{\nano\second} when it in
fact cost \SI{0.53}{\nano\second}---wrong by about five nanoseconds and
in the wrong direction. Timing-driven placement is what helped.

\paragraph{Hold is met in the shipping layout and lost in one variant.}
Repairing that variant's last two antenna violations required
\num{60} additional diodes, whose load breaks hold at
\SI{-0.44}{\nano\second} on \num{55} endpoints. The trade is structural:
a gentler repair setting converges on a byte-identical layout. The
layout reported above is the one that meets hold and carries two antenna
violations; the alternative is reported rather than hidden.

\paragraph{It is not a flight part and it is not radiation-hardened.}
No silicon. No beam. The process kit publishes no total-ionising-dose or
single-event data, so the word ``rad-hard'' is not used of this part and
every reliability number here is a seeded injection on a model.

\paragraph{The spacecraft interfaces are mostly absent.} SpaceWire, CAN,
SPI, I2C and a second UART are unbuilt; four candidates are fetched,
elaborated and priced, and nine of the sixteen peripheral slots are
reserved addresses with nothing behind them.

\paragraph{The accelerator is one node, and the mesh is a
specification.} The inter-node link is frozen as an interface---a
\num{24}-bit link word, a valid/ready handshake, a node identifier and
per-node fault counters---and no RTL implements it. In single-node
builds the routing prefix is tied off at the boundary.

\paragraph{There is no DFT.} No scan, no memory built-in self-test, no
test mode. Both clock gates take their test enable tied low.

\paragraph{One clock, one mode.} No clock gating beyond two instances,
no power domains, no retention, no sleep.

\section{Related work}
\label{sec:related}

Spiking accelerators on open PDKs through open flows are established:
OpenSpike~\cite{openspike2023} taped one out in SkyWater
\SI{130}{\nano\meter} at roughly thirty times this node's scale in 2023,
and further open \SI{130}{\nano\meter} neuromorphic silicon has
followed~\cite{kumaresan2026}. The neuromorphic-for-space framing is a
survey topic~\cite{naoukin2023}. An SNN on an
open PDK is not a contribution and is not claimed as one here.

Neuromorphic hardware under radiation has beam data that this work does
not: an open SNN processor~\cite{frenkel2019} has been characterised in a
neutron beam and cross-sections extracted~\cite{nijsink2026}. That is the tier above this work and the
difference is not one of degree.

Radiation-tolerant neuromorphic silicon exists in fabricated form, and
GR801 is the commercial reference named in \S\ref{sec:intro}.

\paragraph{The nearest neighbour is the author's own.}
ZIRH-2~\cite{zirh2_2026} is a
rad-hard experiment chip on this same open PDK, public under
Apache-2.0, carrying a triplicated SERV RISC-V, TMR chains placed in two
deliberately different arrangements and SECDED RAM with scrub-on-read,
through an open flow to a routed layout. That is nearly every property
this design holds except the spiking accelerator and the management
core's scale. It is cited because a design paper that does not name its
author's adjacent chip invites the reading that it was not looked for.

\paragraph{What is left, then.} Not the architecture, and not the
substrate. What this design offers is a complete, open, re-derivable
record of one fault-tolerance scheme taken from register transfer to
routed layout with its costs measured at each boundary and its failures
kept---including the layout that does not close, the remedies that made
it worse, and the campaign that was blind.

\section{Artefacts}
\label{sec:artefacts}

The design, the harnesses, the flow configurations, the measurement
records and both companion papers are published. Numbers printed here
are either re-derived from the tree by a checked registry or marked as
read by hand from a named artefact; the ratio is printed rather than
implied.

\bibliographystyle{plain}
\bibliography{refs}

\end{document}
